\documentclass[11pt]{article}

\usepackage[T1]{fontenc}
\usepackage{lmodern}
\usepackage{microtype}
\usepackage{amsmath,amssymb,amsthm,mathtools,mathrsfs}
\usepackage[a4paper,margin=30mm]{geometry}
\usepackage{array}
\usepackage{xurl}
\usepackage{xcolor}
\usepackage[colorlinks=true,linkcolor=blue!55!black,
            citecolor=blue!55!black,urlcolor=blue!55!black]{hyperref}
\hypersetup{
  pdftitle={Certified Local Weil Positivity in the Even Pole-Moment-Null Sector through Radius 3/5},
  pdfauthor={Sharan Thota}
}

\newtheorem{theorem}{Theorem}[section]
\newtheorem{proposition}[theorem]{Proposition}
\newtheorem{lemma}[theorem]{Lemma}
\newtheorem{corollary}[theorem]{Corollary}
\theoremstyle{remark}
\newtheorem{remark}[theorem]{Remark}

\DeclareMathOperator{\Tr}{Tr}
\DeclareMathOperator{\sinc}{sinc}
\DeclareMathOperator{\Dom}{Dom}

\setlength{\emergencystretch}{3em}

\begin{document}

\title{Certified Local Weil Positivity in the Even Pole-Moment-Null
Sector through Radius \(3/5\)}
\author{Sharan Thota}
\date{}

\maketitle

\begin{abstract}
We study the compactly supported Weil quadratic form associated with the Riemann zeta function after restriction to even functions whose \(\cosh(x/2)\)-moment vanishes.
For a support radius \(a\), its Fourier multiplier consists of the archimedean digamma term and the prime-power atoms with \(\log n<2a\).
We prove, with a reproducible interval-arithmetic certificate,
that the resulting closed form satisfies
\[
        T_a[f]\geq \frac{3}{100000}\,\lVert f\rVert_2^2
        \qquad (0<a\leq 3/5).
\]
At the endpoint, the arithmetic part contains exactly the atoms \(n=2,3\).
The proof is an infinite-dimensional estimate, not a finite-section Ritz calculation.
We split the form into a compact adverse operator and a positive favorable form,
retain a trace-class part of the favorable form,
and control the discarded directions by a matrix-valued Schur estimate.
A rank-44 interval \(LDL^{*}\) certificate then gives the stated rational gap.
We also prove a general completeness theorem for this positive-rescue scheme,
derive an exact positive-background Poisson-frame representation,
and record the corresponding normalized shorted-graph Loewner inequality.
The result is local in the support radius; the all-radius positivity criterion,
which is equivalent to the Riemann hypothesis, is not asserted.
\end{abstract}

\medskip
\noindent\textbf{2020 Mathematics Subject Classification.} 65G40, 47A07, 47B65, 11M26.

\smallskip
\noindent\textbf{Keywords.} Validated numerics, interval arithmetic,
localized quadratic form, Schur complement, Poisson frame,
Weil's criterion, Riemann zeta function.

\section{Introduction}
\label{pub-sec-introduction}

Let \(\Lambda\) denote the von Mangoldt function and \(\psi=\Gamma'/\Gamma\) the digamma function.
For \(a>0\), the compact-support explicit formula leads to the real multiplier
\begin{equation}
 b_a(t)=
 \operatorname{Re}\psi\!\left(\frac14+\frac{it}{2}\right)-\log\pi
 -2\sum_{1<n<e^{2a}}\frac{\Lambda(n)}{\sqrt n}
       \cos(t\log n).
 \label{pub-eq-symbol-intro}
\end{equation}
The finite sum changes only when \(2a\) crosses the logarithm of a prime power.
This makes fixed-support positivity a natural local problem:
it retains exact arithmetic coefficients while involving only finitely many prime-power translations.

The purpose of this paper is to establish a strict local theorem after the first two arithmetic atoms have entered.
The principal result is the following.
All notation is defined in Section~\ref{pub-sec-weil-form}.

\begin{theorem}[Prime-inclusive local coercivity]
\label{pub-thm-main}
For every \(0<a\leq3/5\), the closed form \(T_a\) on the complex-even,
\(\cosh(x/2)\)-moment-null space satisfies
\begin{equation}
       T_a[f]\geq\frac{3}{100000}\lVert f\rVert_2^2,
       \qquad f\in\mathcal Q_a^0.
\label{pub-eq-main-bound}
\end{equation}
At \(a=3/5\), the prime-power sum in \eqref{pub-eq-symbol-intro} contains exactly \(n=2\) and \(n=3\).
\end{theorem}

The proof has three parts.
First, a projected cosine-frame representation turns multiplication by the negative part of the symbol into a positive trace-class operator.
Second, a favorable, sign-certified portion of the positive part is retained without discarding its orientation relative to the adverse operator.
A matrix Schur estimate controls the infinite orthogonal complement by trace tails.
Third, outward-rounded Arb arithmetic verifies the resulting finite matrix inequalities and every scalar sign decision.
Exact zero-extension then transfers the endpoint bound to all smaller supports,
including across the thresholds \((\log2)/2\) and \((\log3)/2\).

The finite-dimensional space in the certificate is used only to bound the complement of a trace-class operator.
No smallest eigenvalue of a compression of \(T_a\) is used as a lower bound for the continuum form;
such a use would have the wrong variational direction.

The connection with the Riemann hypothesis is contextual rather than a claim of this paper.
Weil's positivity criterion, in its localized Hermitian-form version,
identifies the Riemann hypothesis with positivity for every support radius.
Theorem~\ref{pub-thm-main} supplies one explicit initial interval of radii and therefore does not settle the all-radius problem;
the precise statement is recorded in Section~\ref{pub-sec-conclusion}.

The localization of Weil's quadratic form was initiated by Yoshida,
who proved positivity for sufficiently small support and reduced the problem to finite calculations on suitable localized spaces.
Bombieri subsequently developed the associated variational problem and recorded the classical calculation at the prime-free radius \(a=(\log 2)/2\) \cite{Yoshida1992,Bombieri2000,Bombieri2003}.
Connes and Consani gave an operator-theoretic explanation of the archimedean inequality at that scale and later reported numerical semi-local evidence across successive arithmetic thresholds \cite{ConnesConsani2021,ConnesConsani2023}.
Connes and van Suijlekom studied the quadratic forms attached to finite truncations and proved an extension of the Carath\'eodory--Fej\'er theorem for Toeplitz matrices supplying the required self-adjointness \cite{ConnesVanSuijlekom2025}.
Building on that framework, Connes, Consani and Moscovici construct self-adjoint rank-one perturbations of the scaling spectral triple on an interval whose spectra approach the low-lying zeros;
their convergence is reported numerically, and a rigorous proof of it would give the Riemann hypothesis \cite{ConnesConsaniMoscovici2025}.
Suzuki constructed the localized Friedrichs operator and proved continuity of its lowest eigenvalue,
but his unconditional positivity theorem remains qualitative for sufficiently small \(a\) \cite{Suzuki2023,Suzuki2026}.
A recent finite-Galerkin study supplies an exact dictionary between real even coefficient vectors and band-limited Guinand--Weil test functions, together with a certification rule for finite-cutoff positivity,
but not a lower bound for the exact infinite-dimensional pole-moment-null continuum form \cite{Groskin2026}.

To the author's knowledge, Theorem~\ref{pub-thm-main} is the first explicit rigorous continuum lower bound in the even pole-moment-null sector at a support radius beyond the \(n=3\) entry threshold.
This comparison is restricted to that sector;
it does not concern Suzuki's unrestricted localized operator or the all-radius problem.

The paper is organized as follows.
Section~\ref{pub-sec-weil-form} fixes the Fourier normalization and the form domain.
Section~\ref{pub-sec-cosine-frame} gives the exact projected cosine kernel.
Section~\ref{pub-sec-rescue} proves the abstract positive-rescue estimates.
Section~\ref{pub-sec-certificate} gives the complete theorem-bearing certificate at \(a=3/5\).
Section~\ref{pub-sec-nesting} proves inward nesting,
Section~\ref{pub-sec-poisson-frame} gives the exact grouped Poisson frame,
and Section~\ref{pub-sec-shorting} records the shorted-form consequence.
Section~\ref{pub-sec-conclusion} records the all-radius criterion that delimits the scope of the result.
Reproduction data appear in Appendix~\ref{pub-sec-reproduction}.

\section{The local Weil form}
\label{pub-sec-weil-form}

Write \(I_a=(-a,a)\), and identify every function on \(I_a\) with its zero extension to \(\mathbb R\).
We use
\begin{equation}
 \widehat f(t)=\int_{-a}^{a}f(x)e^{-itx}\,dx,
 \qquad
 \lVert f\rVert_2^2=\frac1{2\pi}
        \int_{\mathbb R}|\widehat f(t)|^2\,dt.
 \label{pub-eq-fourier-normalization}
\end{equation}
Let
\begin{equation}
 \mathcal V_a^0=
 \left\{f\in L^2(I_a;\mathbb C):
 f(-x)=f(x)\ \text{a.e.},\quad
 \int_{-a}^{a}f(x)\cosh(x/2)\,dx=0\right\}.
 \label{pub-eq-moment-space}
\end{equation}
For an even function the odd \(\sinh(x/2)\)-moment vanishes, and hence
\begin{equation}
 \widehat f(i/2)=\widehat f(-i/2)=0
 \quad\Longleftrightarrow\quad
 \int_{-a}^{a}f(x)\cosh(x/2)\,dx=0.
 \label{pub-eq-two-poles-one-moment}
\end{equation}
Thus the single condition in \eqref{pub-eq-moment-space} removes both pole terms from the even restriction of the Weil form.

Put
\begin{equation}
 g_\infty(t)=
 \operatorname{Re}\psi\!\left(\frac14+\frac{it}{2}\right)-\log\pi.
 \label{pub-eq-g-infty}
\end{equation}
Since \(g_\infty(t)=\log(2+|t|)+O(1)\),
while the arithmetic part of \eqref{pub-eq-symbol-intro} is bounded for fixed \(a\),
the natural closed form domain is
\begin{equation}
 \mathcal Q_a^0=
 \left\{f\in\mathcal V_a^0:
 \int_{\mathbb R}\log(2+|t|)|\widehat f(t)|^2\,dt<\infty\right\}.
 \label{pub-eq-form-domain}
\end{equation}
On this domain define
\begin{equation}
 T_a[f,g]=\frac1{2\pi}\int_{\mathbb R}
     b_a(t)\widehat f(t)\overline{\widehat g(t)}\,dt,
 \qquad T_a[f]=T_a[f,f].
 \label{pub-eq-local-form}
\end{equation}
The form is densely defined, closed, and bounded below.
Including an atom with \(\log n=2a\) does not alter the compressed form:
the corresponding translation meets the support only on a null set.

\section{Projected cosine frames}
\label{pub-sec-cosine-frame}

Let \(m_a(x)=\cosh(x/2)\), and let \(P_a^0\) be the orthogonal projection from the even subspace of \(L^2(I_a)\) onto \(m_a^\perp\).
For \(t\geq0\) put
\begin{equation}
       v_{a,t}=P_a^0\cos(tx).
\label{pub-eq-projected-feature}
\end{equation}
For \(f\in\mathcal V_a^0\), evenness and the moment constraint give
\begin{equation}
       \widehat f(t)=\langle f,v_{a,t}\rangle.
\label{pub-eq-analysis-feature}
\end{equation}
Consequently, if \(w\geq0\) is integrable and compactly supported, then
\begin{equation}
 K_{a,w}=\frac1\pi\int_0^\infty
       w(t)|v_{a,t}\rangle\langle v_{a,t}|\,dt
\label{pub-eq-frame-operator}
\end{equation}
is positive and trace class, with
\begin{equation}
 \langle K_{a,w}f,f\rangle=\frac1\pi\int_0^\infty
       w(t)|\widehat f(t)|^2\,dt,
 \qquad
 \Tr K_{a,w}=\frac1\pi\int_0^\infty
       w(t)\lVert v_{a,t}\rVert^2\,dt.
 \label{pub-eq-frame-form-trace}
\end{equation}

The kernel needed by the certificate is elementary.
With \(\sinc z=\sin z/z\), continuously extended at zero,
\begin{align}
 \lVert m_a\rVert_2^2&=a+\sinh a,
 \label{pub-eq-moment-norm}\\
 \langle m_a,\cos(tx)\rangle
 &=2\frac{\frac12\sinh(a/2)\cos(at)
       +t\cosh(a/2)\sin(at)}{t^2+1/4},
 \label{pub-eq-moment-cosine}\\
 \langle v_{a,s},v_{a,t}\rangle
 &=a\{\sinc(a(s-t))+
       \sinc(a(s+t))\}
 \notag\\
 &\hspace{1cm}-
 \frac{\langle m_a,\cos(sx)\rangle
       \langle m_a,\cos(tx)\rangle}{a+\sinh a}.
 \label{pub-eq-projected-kernel}
\end{align}
These identities make both the Gram matrix and all operator compressions amenable to rigorous interval integration.

\section{Positive rescue and a matrix Schur estimate}
\label{pub-sec-rescue}

\subsection{The canonical split}

Let \(b\) be any real multiplier satisfying \(b(t)\to+\infty\), fix \(\mu>0\), and set
\begin{equation}
 w=(\mu-b)_+,
 \qquad d=(b-\mu)_+.
\label{pub-eq-positive-negative-split}
\end{equation}
On a projected cosine space as above, let
\begin{equation}
 A=K_w,
 \qquad
 D[f]=\frac1\pi\int_0^\infty d(t)|\widehat f(t)|^2\,dt.
\label{pub-eq-adverse-favorable}
\end{equation}
Then \(A\) is positive trace class, \(D\) is a closed positive form,
and the multiplier form has the exact decomposition
\begin{equation}
             T=\mu I-A+D.
\label{pub-eq-canonical-rescue-split}
\end{equation}
If \(C\) is positive trace class and \(0\preceq C\preceq D\) in form sense, then
\begin{equation}
             T\succeq\mu I-(A-C).
\label{pub-eq-retained-rescue}
\end{equation}

\subsection{The matrix residual}

We first isolate the operator-theoretic estimate used in the certificate.

\begin{lemma}[Schur residual]
\label{pub-lem-schur-residual}
Let \(S=S^*\) be bounded, let \(P\) be an orthogonal projection, and put \(Q=I-P\).
Write
\[
 D_P=PSP,\qquad X=QSP,\qquad E=QSQ.
\]
Suppose that \(E\preceq\delta Q\) and \(X^*X\preceq R_P\),
where \(R_P\succeq0\) acts on \(P\mathcal H\).
If \(\lambda>\delta\) and
\begin{equation}
 \lambda P-D_P-\frac{R_P}{\lambda-\delta}\succeq0,
\label{pub-eq-schur-residual-condition}
\end{equation}
then \(S\preceq\lambda I\).
\end{lemma}

\begin{proof}
Relative to \(P\mathcal H\oplus Q\mathcal H\),
\[
 \lambda I-S=
 \begin{pmatrix}
   \lambda P-D_P&-X^*\\
   -X&\lambda Q-E
 \end{pmatrix}.
\]
The lower-right block is at least \((\lambda-\delta)Q\).
Its inverse is therefore at most \((\lambda-\delta)^{-1}Q\),
and its Schur complement is bounded below by the left side of \eqref{pub-eq-schur-residual-condition}.
The Schur complement criterion proves the assertion.
\end{proof}

\begin{theorem}[Positive-pair matrix residual]
\label{pub-thm-positive-pair}
Let \(A,C\succeq0\) be trace-class operators, set \(S=A-C\),
and let \(P+Q=I\) be an orthogonal decomposition.
Put
\begin{equation}
 \tau_A=\Tr(QAQ),
 \qquad \tau_C=\Tr(QCQ).
\label{pub-eq-trace-tails}
\end{equation}
For \(\eta>0\), define on \(P\mathcal H\)
\begin{equation}
 R_\eta=(1+\eta)\tau_A PAP
      +(1+\eta^{-1})\tau_C PCP.
\label{pub-eq-matrix-residual}
\end{equation}
If \(\lambda>\tau_A\) and
\begin{equation}
 \lambda P-P(A-C)P-\frac{R_\eta}{\lambda-\tau_A}\succeq0,
\label{pub-eq-positive-pair-lmi}
\end{equation}
then \(A-C\preceq\lambda I\).
\end{theorem}

\begin{proof}
For every positive block operator \(B\), Douglas factorization \cite{Douglas1966} gives
\begin{equation}
 (QBP)^*(QBP)\preceq\lVert QBQ\rVert PBP
 \preceq\Tr(QBQ)PBP.
\label{pub-eq-positive-block-factorization}
\end{equation}
Apply this to \(A\) and \(C\).
Matrix Young's inequality gives
\[
 [Q(A-C)P]^*[Q(A-C)P]
 \preceq(1+\eta)(QAP)^*(QAP)
 +(1+\eta^{-1})(QCP)^*(QCP)
 \preceq R_\eta.
\]
Moreover, \(Q(A-C)Q\preceq QAQ\preceq\tau_AQ\).
The conclusion follows from Lemma~\ref{pub-lem-schur-residual}.
\end{proof}

\subsection{Completeness at a fixed support}

The next theorem explains why bounded positive rescues and finite projectors do not create an obstruction at a fixed radius where strict coercivity already holds.
It is a completeness statement for the proof scheme,
not an independent positivity theorem.

\begin{theorem}[Exhaustion of an unbounded favorable form]
\label{pub-thm-rescue-completeness}
Let \(A\succeq0\) be compact, let \(D\) be a closed positive form,
and let \(C_R\succeq0\) be compact operators increasing monotonically to \(D\) in form sense.
Define
\begin{equation}
 s_*=\sup_{\substack{f\in\Dom D^{1/2}\\\lVert f\rVert=1}}
       \bigl(\langle Af,f\rangle-D[f]\bigr),
 \qquad
 \lambda_R=\sup\sigma(A-C_R).
\label{pub-eq-completeness-quantities}
\end{equation}
If the Hilbert space is infinite dimensional, then
\begin{equation}
                \lambda_R\downarrow\max(0,s_*).
\label{pub-eq-completeness-limit}
\end{equation}
In particular, if \(\mu I-A+D\succeq cI\) for some \(\mu,c>0\),
then \(\lambda_R<\mu\) for all sufficiently large \(R\).
\end{theorem}

\begin{proof}
Monotonicity gives \(\lambda_R\downarrow\ell\geq0\),
because \(A-C_R\) is compact and hence \(0\) lies in its spectrum.
The form inequality \(C_R\preceq D\) gives \(\ell\geq\max(0,s_*)\).

Suppose \(\ell>0\), choose \(R_n\to\infty\), and let \(f_n\) be unit top eigenvectors.
After passage to a subsequence, \(f_n\rightharpoonup f\).
Since \(C_{R_n}\succeq0\), \(\langle Af_n,f_n\rangle\geq\lambda_{R_n}\to\ell\).
Compactness of \(A\) implies \(f\neq0\).
For fixed \(R_k\) and sufficiently large \(n\),
\[
 \lambda_{R_n}\leq
 \langle Af_n,f_n\rangle-\langle C_{R_k}f_n,f_n\rangle.
\]
Compactness of \(A\) and \(C_{R_k}\), followed by monotone form convergence in \(k\), yields
\[
 \ell\leq\langle Af,f\rangle-D[f]
 \leq s_*\lVert f\rVert^2\leq s_*.
\]
Thus \(\ell=s_*>0\).
When \(\ell=0\), the preceding lower bound already gives \(\max(0,s_*)=0\),
proving \eqref{pub-eq-completeness-limit}.
Finally, strict coercivity gives \(s_*\leq\mu-c\), so the limit is strictly below \(\mu\).
\end{proof}

\begin{corollary}[Finite-projector completeness]
\label{pub-cor-finite-projector-completeness}
Under the final hypothesis of Theorem~\ref{pub-thm-rescue-completeness},
assume moreover that \(A\) and the selected \(C_R\) are trace class,
and choose \(R\) so that \(\lambda_R<\mu\).
If \(P_N\uparrow I\) are finite-rank orthogonal projections,
then the trace tails in Theorem~\ref{pub-thm-positive-pair} tend to zero and its finite LMI certifies a bound strictly below \(\mu\) for all sufficiently large \(N\).
\end{corollary}

\begin{proof}
For the trace-class operators \(A\) and \(C_R\),
\(\Tr((I-P_N)A(I-P_N))\to0\) and likewise for \(C_R\).
Since \(A-C_R\) is compact, \(P_N(A-C_R)P_N\to A-C_R\) in norm.
The residual in \eqref{pub-eq-matrix-residual} therefore tends to zero,
and the finite Schur test converges to the strict inequality \(\lambda_R<\mu\).
\end{proof}

\section{The certified endpoint \texorpdfstring{\(a=3/5\)}{a=3/5}}
\label{pub-sec-certificate}

\subsection{The exact arithmetic cell}

Set \(a_0=3/5\).
Rigorous interval evaluation gives
\begin{equation}
             \log3<\frac65<\log4.
\label{pub-eq-cell-inequalities}
\end{equation}
Thus the endpoint symbol is
\begin{equation}
 b_{2,3}(t)=g_\infty(t)
 -\sqrt2\log2\cos(t\log2)
 -\frac{2\log3}{\sqrt3}\cos(t\log3).
\label{pub-eq-b23}
\end{equation}
Put \(d=1-b_{2,3}\).
In the notation of \eqref{pub-eq-canonical-rescue-split},
\begin{equation}
 T_{a_0}=I-A+C_{\mathrm{full}},
 \qquad A=K_{a_0,d_+},
 \qquad C_{\mathrm{full}}=K_{a_0,(-d)_+}.
\label{pub-eq-endpoint-full-split}
\end{equation}

The retained rescue \(C\preceq C_{\mathrm{full}}\) is obtained by restricting \((-d)_+\) to the following union \(J\) of exact rational intervals:
\begin{equation}
\begin{split}
J={}&[60382193/20000000,\,380082633/100000000]\\
&\cup[39235217/3125000,\,1573243267/100000000]\\
&\cup[1908657941/100000000,\,228037531/10000000]\\
&\cup[23057517/1000000,\,1340114843/50000000]\\
&\cup[184766397/6250000,\,3362720043/100000000]\\
&\cup[717875089/20000000,\,4443581909/100000000].
\end{split}
\label{pub-eq-rescue-cores}
\end{equation}
Interval subdivision proves \(d<0\) throughout each component of \(J\), so
\begin{equation}
 C=K_{a_0,(-d)\mathbf1_J}
 \qquad\text{satisfies}\qquad 0\preceq C\preceq C_{\mathrm{full}}.
\label{pub-eq-retained-C}
\end{equation}

\subsection{Complete sign and tail control}

The derivative of the archimedean term has the positive expansion
\begin{equation}
 g_\infty'(t)=\frac t2\sum_{k\geq0}
 \frac{k+1/4}{((k+1/4)^2+(t/2)^2)^2}>0
 \qquad(t>0).
\label{pub-eq-g-derivative}
\end{equation}
Using
\[
 \sup_{y\geq0}\frac{yu}{(u^2+y^2)^2}
 =\frac{9}{16\sqrt3}\frac1{u^2},
 \qquad
 \sum_{k\geq0}\frac1{(k+1/4)^2}=\pi^2+8G,
\]
where \(G\) is Catalan's constant, one obtains the global rationally certified estimate
\begin{equation}
 \lVert d'\rVert_\infty
 \leq\frac{9(\pi^2+8G)}{16\sqrt3}
 +\sqrt2(\log2)^2+\frac{2(\log3)^2}{\sqrt3}
 <7.658125<8.
\label{pub-eq-lipschitz-eight}
\end{equation}

Outward-rounded interval subdivision on \([0,256]\) produces 27 disjoint rational root brackets,
each of width \(10^{-10}\), with alternating strict endpoint signs,
and proves the sign of \(d\) on all 28 complementary intervals.
The exact list of brackets and the complete subdivision tree are part of the supplementary proof ledger described in Appendix~\ref{pub-sec-reproduction}.
Formula \eqref{pub-eq-lipschitz-eight} and \(\lVert v_{a_0,t}\rVert^2\leq2a_0=6/5\) bound the total positive operator possibly omitted inside all root brackets by
\begin{equation}
 0\preceq R_{\mathrm{root}}\preceq
 \varepsilon_{\mathrm{root}}I,
 \qquad
 \varepsilon_{\mathrm{root}}<
 \frac{27\cdot8\cdot(6/5)\cdot10^{-20}}{\pi}
 <8.250593\times10^{-19}.
\label{pub-eq-root-charge}
\end{equation}

Finally, monotonicity of \(g_\infty\) and the elementary bound \(|\cos|\leq1\) give
\begin{equation}
 1-g_\infty(256)+\sqrt2\log2+\frac{2\log3}{\sqrt3}
 <-0.4584733976219665.
\label{pub-eq-tail-sign}
\end{equation}
Hence \(d(t)<0\) for every \(t\geq256\), and the adverse operator has no unaccounted tail.

\subsection{The rank-44 continuum projection}

Let
\begin{equation}
 \xi_j=\frac{(j+1/2)\pi}{a_0}
      =\frac{5(2j+1)\pi}{6},
 \qquad 0\leq j<44,
\label{pub-eq-feature-frequencies}
\end{equation}
and define \(V:\mathbb C^{44}\to\mathcal V_{a_0}^0\) by \(Ve_j=v_{a_0,\xi_j}\).
Put
\begin{equation}
 G=V^*V,
 \qquad P=VG^{-1}V^*,
 \qquad Q=I-P.
\label{pub-eq-proof-projection}
\end{equation}
The unprojected cosines form an exact orthogonal half-Dirichlet family;
the moment projection produces the explicitly computable Gram matrix \(G\).
Interval \(LDL^*\) proves \(G\succ0\), with smallest pivot greater than
\begin{equation}
                   0.11996777343402634.
\label{pub-eq-gram-pivot}
\end{equation}
For a positive trace-class operator \(B\), if \(B_P=V^*BV\), then
\begin{equation}
              \Tr(PBP)=
              \Tr(G^{-1}B_P).
\label{pub-eq-captured-trace}
\end{equation}
This identity converts outward enclosures of full traces and inward enclosures of captured traces into rigorous upper bounds for the infinite-dimensional complement.

Let \(A_0\) denote the adverse operator obtained from all sign-certified positive intervals,
excluding only the root brackets, and write \(B_{A,0}=V^*A_0V\), \(B_C=V^*CV\).
The interval integrations and trace subtractions give
\begin{equation}
\begin{aligned}
 \Tr A&<4.170131812888502405,\\
 \Tr(PAP)&>4.170127556767087307,\\
 \tau_A=\Tr(QAQ)
   &<4.256121415098\times10^{-6}<1/200000,\\[1mm]
 \Tr C&<5.040647295020747,\\
 \Tr(PCP)&>5.040645837226345324,\\
 \tau_C=\Tr(QCQ)
   &<1.457794400983\times10^{-6}<1/500000.
\end{aligned}
\label{pub-eq-certified-traces}
\end{equation}
All subtractions in \eqref{pub-eq-certified-traces} are one-sided:
the full trace is rounded upward and the captured trace downward.
The root charge \eqref{pub-eq-root-charge} is added only in the adverse direction.

For
\begin{equation}
 \eta=\frac1{16},
 \qquad
 \lambda=\frac{99997}{100000},
 \qquad
 \bar\tau_A=\frac1{200000},
 \qquad
 \bar\tau_C=\frac1{500000},
\label{pub-eq-lmi-rationals}
\end{equation}
define the Gram-coordinate matrix
\begin{align}
 \mathcal M={}&
 \lambda G-(B_{A,0}+\varepsilon_{\mathrm{root}}G-B_C)
 \notag\\
 &-\frac{
 (1+\eta)\bar\tau_A(B_{A,0}+\varepsilon_{\mathrm{root}}G)
 +(1+\eta^{-1})\bar\tau_CB_C}
 {\lambda-\bar\tau_A}.
\label{pub-eq-final-lmi}
\end{align}
An outward-rounded interval \(LDL^*\) factorization gives 44 strictly positive pivots,
the smallest with lower endpoint greater than
\begin{equation}
                  0.00979334885665161.
\label{pub-eq-final-pivot}
\end{equation}

\begin{proposition}[Endpoint certificate]
\label{pub-prop-endpoint-certificate}
On \(\mathcal Q_{3/5}^0\),
\begin{equation}
                 T_{3/5}\succeq\frac3{100000}I.
\label{pub-eq-endpoint-bound}
\end{equation}
\end{proposition}

\begin{proof}
The sign proof gives \(A\preceq A_0+R_{\mathrm{root}}\),
while \eqref{pub-eq-retained-C} gives \(C\preceq C_{\mathrm{full}}\).
Thus
\[
 T_{3/5}\succeq I-(A-C).
\]
Equations \eqref{pub-eq-certified-traces}--\eqref{pub-eq-final-pivot} verify the hypothesis of Theorem~\ref{pub-thm-positive-pair}, so
\[
 A-C\preceq\frac{99997}{100000}I.
\]
Subtracting this inequality from the identity proves \eqref{pub-eq-endpoint-bound}.
\end{proof}

\section{Exact inward nesting}
\label{pub-sec-nesting}

The passage from the endpoint to all smaller supports is exact and does not use numerical continuity in \(a\).

\begin{lemma}[Zero-extension identity]
\label{pub-lem-zero-extension}
Let \(0<a\leq a_0\), and let \(f\in\mathcal Q_a^0\).
If \(f^{\mathrm{ext}}\) denotes its zero extension from \((-a,a)\) to \((-a_0,a_0)\), then
\begin{equation}
 \lVert f^{\mathrm{ext}}\rVert_2=\lVert f\rVert_2,
 \qquad
 T_{a_0}[f^{\mathrm{ext}}]=T_a[f].
\label{pub-eq-zero-extension-identity}
\end{equation}
\end{lemma}

\begin{proof}
Zero extension preserves evenness, the \(\cosh(x/2)\)-moment, the norm, and the Fourier transform.
It remains to compare the arithmetic sums.
If an atom \(n\) is active at \(a_0\) but inactive at \(a\), then \(\log n\geq2a\).
Its contribution to the form is a constant multiple of the correlation of \(f\) with its translate by \(\log n\).
The two supports are disjoint, except possibly at a null endpoint when \(\log n=2a\),
so this correlation is zero.
All other terms agree.
\end{proof}

\begin{proof}[Proof of Theorem~\ref{pub-thm-main}]
Apply Proposition~\ref{pub-prop-endpoint-certificate} to \(f^{\mathrm{ext}}\) with \(a_0=3/5\),
and use Lemma~\ref{pub-lem-zero-extension}.
\end{proof}

\section{An exact positive-background Poisson frame}
\label{pub-sec-poisson-frame}

The prime-power part of the local symbol can be grouped by primes without losing any coefficient.
This gives an exact difference of two positive forms and makes the shorted-form consequence in the next section transparent.

Let \(\iota_a:L^2(I_a)\to L^2(\mathbb R)\) be zero extension, let \(U_hf(x)=f(x-h)\), and put
\begin{equation}
 X=e^{2a},\qquad h_p=\log p,\qquad r_p=p^{-1/2},\qquad
 \vartheta(X)=\sum_{p\leq X}\log p.
\label{pub-eq-poisson-parameters}
\end{equation}
Define the compressed Poisson Gramian
\begin{equation}
 \Pi_{p,a}=\iota_a^*(1-r_p^2)
 (I-r_pU_{h_p})^{-*}(I-r_pU_{h_p})^{-1}\iota_a.
\label{pub-eq-poisson-gramian}
\end{equation}
The archimedean positive form is
\begin{equation}
 \mathcal J_{\infty,a}[f]=
 \int_0^\infty\frac{e^{-h/2}}{1-e^{-2h}}
 \lVert(I-U_h)\iota_af\rVert_2^2\,dh.
\label{pub-eq-archimedean-jump}
\end{equation}
It is understood as a closed form on the logarithmic domain \eqref{pub-eq-form-domain}.
Put
\begin{equation}
 \kappa_a=\vartheta(X)+g_\infty(0),\qquad
 s\geq\max(0,\kappa_a),\qquad
 \delta_s=s-\kappa_a,
\label{pub-eq-background-shift}
\end{equation}
and set
\begin{equation}
 R_{s,a}=\mathcal J_{\infty,a}+sI,
 \qquad
 \mathscr A_{X,s,a}=\delta_sI+
       \sum_{p\leq X}(\log p)\Pi_{p,a}.
\label{pub-eq-positive-background-forms}
\end{equation}
Both forms in \eqref{pub-eq-positive-background-forms} are positive.

\begin{proposition}[Coefficient-locked Poisson-frame identity]
\label{pub-prop-poisson-frame}
On \(\mathcal Q_a^0\),
\begin{equation}
                    T_a=R_{s,a}-\mathscr A_{X,s,a}.
\label{pub-eq-exact-poisson-split}
\end{equation}
Moreover, \(\mathscr A_{X,s,a}=F_{s,a}^*F_{s,a}\), where
\begin{equation}
 F_{s,a}f=\left(\sqrt{\delta_s}\,f,
 \left(\sqrt{(\log p)(1-r_p^2)}
 (I-r_pU_{h_p})^{-1}\iota_af\right)_{p\leq X}\right).
\label{pub-eq-poisson-analysis-operator}
\end{equation}
Thus \eqref{pub-eq-exact-poisson-split} is an exact positive-background frame representation, not a primewise estimate.
\end{proposition}

\begin{proof}
The Neumann series of \((I-r_pU_{h_p})^{-1}\) gives the Poisson-kernel identity
\begin{equation}
 \langle f,\Pi_{p,a}f\rangle
 =\lVert f\rVert_2^2+2\sum_{k\geq1}p^{-k/2}
   \operatorname{Re}\langle\iota_af,U_{kh_p}\iota_af\rangle.
\label{pub-eq-grouped-prime-tower}
\end{equation}
When \(kh_p\geq2a\), the correlation on the right vanishes by disjoint support.
Multiplication by \(\log p\), followed by summation over \(p\leq X\),
therefore reproduces exactly the von Mangoldt coefficients of every active prime power in \eqref{pub-eq-symbol-intro}.

The integral representation of the digamma function and Parseval's identity give
\begin{equation}
 \frac1{2\pi}\int_{\mathbb R}g_\infty(t)|\widehat f(t)|^2\,dt
 =\mathcal J_{\infty,a}[f]+g_\infty(0)\lVert f\rVert_2^2.
\label{pub-eq-digamma-jump-identity}
\end{equation}
Combining \eqref{pub-eq-grouped-prime-tower} and \eqref{pub-eq-digamma-jump-identity},
and using \(s-\delta_s=\vartheta(X)+g_\infty(0)\),
proves \eqref{pub-eq-exact-poisson-split}.
Formula \eqref{pub-eq-poisson-analysis-operator} directly gives the stated Gram factorization.
\end{proof}

\begin{remark}
Expanding one prime tower and truncating it before compression would introduce a remainder.
Equation \eqref{pub-eq-grouped-prime-tower} instead groups the entire geometric tower first;
compact support then annihilates every inactive power exactly.
\end{remark}

\section{A normalized shorted-form consequence}
\label{pub-sec-shorting}

For completeness we record the abstract shorted-form consequence that motivates the normalized Loewner formulation.
It requires no additional arithmetic estimate.
All spaces and maps in this section are restricted to the certified even space \(\mathcal V_a^0\);
no assertion is made for odd directions in the full two-pole space.

Let \(T\) be a positive closed form on a Hilbert space \(\mathcal H\),
with form domain \(\mathcal Q\), and suppose \(T\succeq cI\).
Let \(P\) be a finite-rank orthogonal projection such that \(P\mathcal Q\subset\mathcal Q\), and put \(E=I-P\).
For each \(p\in P\mathcal H\),
Lax--Milgram gives a unique \(e_p\in E\mathcal Q\) such that
\begin{equation}
 T[p+e_p,e]=0\qquad(e\in E\mathcal Q).
\label{pub-eq-minimizing-graph-variational}
\end{equation}
Define \(J_Tp=p+e_p\).
Thus \(PJ_T=I_{P\mathcal H}\), and \(J_Tp\) is the unique form-minimizing lift of \(p\).
When the relevant block operators are bounded,
this variational definition agrees with the familiar shorthand
\begin{equation}
 J_T=P-E(ETE)^{-1}ETP.
\label{pub-eq-minimizing-graph}
\end{equation}
Define
\begin{equation}
 B_G=J_T^*J_T\succeq I,
 \qquad
 \Gamma=J_TB_G^{-1/2}.
\label{pub-eq-normalized-graph}
\end{equation}
Then \(\Gamma^*\Gamma=I\).

\begin{corollary}[Normalized shorted-graph Poisson inequality]
\label{pub-cor-normalized-graph}
Suppose in addition that \(T=R-\mathscr A\) as closed forms.
Then
\begin{equation}
          \Gamma^*\mathscr A\Gamma+cI
          \preceq\Gamma^*R\Gamma.
\label{pub-eq-normalized-loewner}
\end{equation}
For the local Weil form and the positive-background decomposition of Proposition~\ref{pub-prop-poisson-frame},
one may take \(R=R_{s,a}\), \(\mathscr A=\mathscr A_{X,s,a}\),
and \(c=3/100000\) for every \(0<a\leq3/5\), with all maps defined on \(\mathcal V_a^0\),
and with \(s\geq\max\{0,\vartheta(e^{2a})+g_\infty(0)\}\).
In particular,
\begin{equation}
 \Gamma_a^*\mathscr A_{X,s,a}\Gamma_a+\frac3{100000}I
 \preceq\Gamma_a^*R_{s,a}\Gamma_a.
\label{pub-eq-specific-poisson-loewner}
\end{equation}
\end{corollary}

\begin{proof}
Coercivity gives \(J_T^*TJ_T\succeq cJ_T^*J_T\).
Congruence by \(B_G^{-1/2}\) yields \(\Gamma^*T\Gamma\succeq cI\).
Substitute \(T=R-\mathscr A\) and rearrange.
This is the usual form-short construction; see \cite{AndersonTrapp1975,Kato1995}.
\end{proof}

\section{Concluding remarks}
\label{pub-sec-conclusion}

Theorem~\ref{pub-thm-main} establishes strict positivity of the exact prime-inclusive Weil form through \(a=3/5\).
The proof separates the analytic issue---control of an infinite-dimensional complement---from the finite interval computation.
Theorem~\ref{pub-thm-rescue-completeness} also shows that,
at a fixed support where strict coercivity is true,
a sufficiently large bounded rescue and finite projector can in principle detect it.

Neither result is uniform in the support radius.
Extending the theorem to all \(a>0\) would require new arithmetic control as further prime-power atoms enter and the low-energy subspace grows.
The following statement, classical in substance and recorded here only to
delimit the scope of Theorem~\ref{pub-thm-main},
identifies that extension as an RH-level problem.

\begin{proposition}[Even local Weil criterion]
\label{pub-prop-rh-equivalence}
The Riemann hypothesis is equivalent to
\begin{equation}
 T_a[f]\geq0
 \quad\text{for every }a>0\text{ and every }f\in\mathcal Q_a^0.
 \label{pub-eq-rh-all-radius}
\end{equation}
It is enough to quantify over smooth functions compactly supported in \(I_a\) and satisfying the moment condition.
\end{proposition}

\begin{proof}
In the compact-support explicit formula of Weil and Yoshida \cite{Weil1952,Yoshida1992},
the non-pole Fourier multiplier is exactly \eqref{pub-eq-symbol-intro}.
Equation \eqref{pub-eq-two-poles-one-moment} removes the two pole evaluations.
Under the Riemann hypothesis,
the zero side of the Hermitian form is a sum of nonnegative squares,
which proves \eqref{pub-eq-rh-all-radius}.

Conversely, the localized interpolation theorem for the Weil form allows one to prescribe the pole zeros and to isolate an off-critical zero orbit.
After conjugation and even symmetrization,
an off-line quartet yields a negative real-even test,
while the other zero evaluations can be made arbitrarily small.
The resulting Paley--Wiener function has some finite exponential type and therefore comes from a test supported in \((-a,a)\) for a finite \(a\).
This contradicts \eqref{pub-eq-rh-all-radius}.
Density of the smooth moment-null core in \(\mathcal Q_a^0\) gives the closed-form statement.
See \cite{Yoshida1992,Bombieri2000} for the localized Hermitian form and interpolation formulation.
\end{proof}

\begin{remark}
Proposition~\ref{pub-prop-rh-equivalence} quantifies over all \(a>0\).
No converse from positivity on a single bounded interval of radii is used or
claimed, and no assertion about the all-radius problem is made in this paper.
\end{remark}

\appendix

\section{Reproduction of the interval certificate}
\label{pub-sec-reproduction}

The theorem-bearing computation uses Python~3.11.5.
Its only external dependency is \texttt{python-flint}~0.9.0, linked against FLINT~3.6.0.
All Arb ball operations use 256-bit working precision.
The generator uses integration tolerance \(10^{-14}\), an evaluation limit of \(500000\),
rank \(44\), and exact rational panels of width at most \(1/2\).
It records every scalar interval, Gram and compression entry, trace enclosure,
sign subdivision, and \(LDL^*\) pivot in a machine-readable JSON ledger.
Arb and its midpoint-radius arithmetic are described in \cite{Johansson2017Arb}.

\paragraph{Code and data availability.}
The paper source, certificate generator, frozen interval ledger,
and separate validator are available in the \href{https://github.com/sharantresearch-ops/certified-local-weil-positivity}{companion repository}.
The version corresponding to this paper is archived under \href{https://github.com/sharantresearch-ops/certified-local-weil-positivity/tree/v1.0.1}{tag \texttt{v1.0.1}}.
The hashes below identify the exact theorem-bearing artifacts.

The supplementary files used for Proposition~\ref{pub-prop-endpoint-certificate} are listed below.
Each hash is SHA--256; line breaks are typographical only.
\begin{description}
\item[\texttt{generator}] \path{certificate/certify.py}\\
{\scriptsize\ttfamily
D9C518F4CE80037E1E545AD7B4282811\allowbreak
9D06A5A53C4F65383A1877E4F7D3BE7D}
\item[\texttt{ledger}] \path{certificate/ledger.json}\\
{\scriptsize\ttfamily
03E81FCA545D184ED320B79E6576ECA7\allowbreak
73976C274C134FBA0C0859D170E4A4B4}
\item[\texttt{validator}] \path{certificate/validate.py}\\
{\scriptsize\ttfamily
8CA47FC4CE1B603756704888F9BCEB41\allowbreak
C89DC112585B25326352DFA0BD65756D}
\item[\texttt{result}] \path{certificate/schur-lmi.json}\\
{\scriptsize\ttfamily
BCC3015588E4E5BE9CE34A362D5CCB56\allowbreak
959B82C12AD1E62381C5F995A2127C10}
\end{description}
The postprocessor verifies the canonical theorem-bearing proof payload
\[
\text{\scriptsize\ttfamily
950A5E6D261E548A40AB776D15F67741\allowbreak
660A9786614E6B000B0CD493378A2285}
\]
before replaying the final serialized interval-matrix LMI.
The generator certifies the root signs, scalar and matrix integrals, Gram positivity,
and one-sided trace bounds recorded in the ledger.
The 60-digit interval balls and the ledger hash provide the data required for independent replay.

After cloning the repository and changing to its top-level directory,
the certificate is reproduced on a POSIX shell by
\begin{verbatim}
python -m pip install -r requirements.txt
mkdir -p build
python certificate/certify.py \
  --workers 8 --precision-bits 256 --tolerance 1e-14 \
  --eval-limit 500000 --json build/fresh-ledger.json
python certificate/validate.py \
  --ledger build/fresh-ledger.json \
  --output build/fresh-schur-lmi.json --precision 256
\end{verbatim}
Parallel worker scheduling affects elapsed time but not the canonical proof payload.

\section{Exact data flow for the continuum bound}
\label{pub-sec-data-flow}

For reference, the logical dependence of the computer-assisted part is
\[
\begin{array}{c}
\text{interval proof of \eqref{pub-eq-cell-inequalities}, signs, and tail}
\\[1mm]\Downarrow\\[-1mm]
A\preceq A_0+\varepsilon_{\mathrm{root}}I,
\quad 0\preceq C\preceq C_{\mathrm{full}}
\\[1mm]\Downarrow\\[-1mm]
\text{interval Gram/compression integrals and one-sided trace tails}
\\[1mm]\Downarrow\\[-1mm]
\mathcal M\succ0\ \text{by interval }LDL^*
\\[1mm]\Downarrow\\[-1mm]
A-C\preceq(99997/100000)I
\\[1mm]\Downarrow\\[-1mm]
T_{3/5}\succeq(3/100000)I
\\[1mm]\Downarrow\\[-1mm]
T_a\succeq(3/100000)I\quad(0<a\leq3/5).
\end{array}
\]
Every arrow except the three interval-verification steps is proved symbolically in the body of the paper.

\begin{thebibliography}{99}

\bibitem{Weil1952}
A.~Weil,
\newblock Sur les ``formules explicites'' de la th\'eorie des nombres
premiers,
\newblock in \emph{Comm. S\'em. Math. Univ. Lund, Tome Suppl\'ementaire},
1952, pp.~252--265.
\newblock \url{https://cds.cern.ch/record/471308}.

\bibitem{Weil1972}
A.~Weil,
\newblock Sur les formules explicites de th\'eorie des nombres,
\newblock \emph{Math. USSR-Izv.} \textbf{6} (1972), 1--17.
\newblock \url{https://doi.org/10.1070/IM1972v006n01ABEH001866}.

\bibitem{Yoshida1992}
H.~Yoshida,
\newblock On Hermitian forms attached to zeta functions,
\newblock in \emph{Zeta Functions in Geometry}, Adv. Stud. Pure Math.
\textbf{21}, Kinokuniya, Tokyo, 1992, pp.~281--325.
\newblock \url{https://doi.org/10.2969/aspm/02110281}.

\bibitem{Bombieri2000}
E.~Bombieri,
\newblock Remarks on Weil's quadratic functional in the theory of prime
numbers, I,
\newblock \emph{Rend. Lincei Mat. Appl.} \textbf{11} (2000), 183--233.
\newblock \url{https://www.bdim.eu/item?id=RLIN_2000_9_11_3_183_0}.

\bibitem{Bombieri2003}
E.~Bombieri,
\newblock A variational approach to the explicit formula,
\newblock \emph{Comm. Pure Appl. Math.} \textbf{56} (2003), 1151--1164.
\newblock \url{https://doi.org/10.1002/cpa.10089}.

\bibitem{ConnesConsani2021}
A.~Connes and C.~Consani,
\newblock Weil positivity and trace formula, the archimedean place,
\newblock \emph{Selecta Math. (N.S.)} \textbf{27} (2021), article 77.
\newblock \url{https://doi.org/10.1007/s00029-021-00689-4}.

\bibitem{ConnesConsani2023}
A.~Connes and C.~Consani,
\newblock Spectral triples and \(\zeta\)-cycles,
\newblock \emph{Enseign. Math.} \textbf{69} (2023), 93--148.
\newblock \url{https://doi.org/10.4171/LEM/1049}.

\bibitem{ConnesVanSuijlekom2025}
A.~Connes and W.~D. van Suijlekom,
\newblock Quadratic forms, real zeros and echoes of the spectral action,
\newblock \emph{Comm. Math. Phys.} \textbf{406} (2025), article 312.
\newblock \url{https://arxiv.org/abs/2511.23257}.

\bibitem{ConnesConsaniMoscovici2025}
A.~Connes, C.~Consani and H.~Moscovici,
\newblock Zeta spectral triples,
\newblock arXiv:2511.22755 (2025); to appear in EMS Lecture Notes in
Mathematics.
\newblock \url{https://arxiv.org/abs/2511.22755}.

\bibitem{Suzuki2023}
M.~Suzuki,
\newblock Aspects of the screw function corresponding to the Riemann
zeta-function,
\newblock \emph{J. Lond. Math. Soc.} \textbf{108} (2023), 1448--1487.
\newblock \url{https://doi.org/10.1112/jlms.12785}.

\bibitem{Suzuki2026}
M.~Suzuki,
\newblock Weil's quadratic form via the screw function,
\newblock arXiv:2606.09096 (2026).
\newblock \url{https://arxiv.org/abs/2606.09096}.

\bibitem{Groskin2026}
A.~Groskin,
\newblock A finite Guinand--Weil dictionary and archimedean tail order for
the truncated Weil quadratic form,
\newblock arXiv:2607.02828 (2026).
\newblock \url{https://arxiv.org/abs/2607.02828}.

\bibitem{Douglas1966}
R.~G. Douglas,
\newblock On majorization, factorization, and range inclusion of operators
on Hilbert space,
\newblock \emph{Proc. Amer. Math. Soc.} \textbf{17} (1966), 413--415.
\newblock \url{https://doi.org/10.1090/S0002-9939-1966-0203464-1}.

\bibitem{AndersonTrapp1975}
W.~N. Anderson, Jr. and G.~E. Trapp,
\newblock Shorted operators II,
\newblock \emph{SIAM J. Appl. Math.} \textbf{28} (1975), 60--71.
\newblock \url{https://doi.org/10.1137/0128007}.

\bibitem{Kato1995}
T.~Kato,
\newblock \emph{Perturbation Theory for Linear Operators},
\newblock Classics in Mathematics, Springer, Berlin, 1995.
\newblock \url{https://doi.org/10.1007/978-3-642-66282-9}.

\bibitem{Johansson2017Arb}
F.~Johansson,
\newblock Arb: efficient arbitrary-precision midpoint-radius interval
arithmetic,
\newblock \emph{IEEE Trans. Comput.} \textbf{66} (2017), 1281--1292.
\newblock \url{https://doi.org/10.1109/TC.2017.2690633}.

\end{thebibliography}

\end{document}
